% GENERATED FILE. Edit canonical lessons, then run npm run generate:books.
% canonical-source-hash: fnv1a64:b1f8d8450d60a4ce
% canonical-lessons: TA-C457-karunai, TA-C457-parivu, TA-C457-anutapam, TA-C457-paritapam, TA-C457-akkarai

\chapter{Mercy and Concern}
\label{ch:ta-mercy-and-concern}
\hlchaptermodality{457}

\begin{chapteropening}
\textbf{By the end of this chapter:} \emph{I can talk about compassion, sympathy and concern for others.}

Complete the last lesson of chapter 457: I can talk about compassion, sympathy and concern for others.
\end{chapteropening}

\section[\texorpdfstring{\ta{கருணை} (karuṇai)}{கருணை (karuṇai)}]{\ta{கருணை} (karuṇai) — mercy; compassion}

\label{lesson:TA-C457-karunai}



\begin{quote}
Before the new one: say the Tamil for consolation, then the Tamil for to subside.
\end{quote}

\subsection*{You'll want to know: \ta{கருணை}}
\textbf{\ta{கருணை}} — \emph{karuṇai} — \textquotedblleft{}mercy; compassion\textquotedblright{}.

\begin{quote}
\ta{அந்த} \ta{மருத்துவர்} \ta{ஏழைகளிடம்} \ta{கருணையுடன்} \ta{நடந்துகொண்டார்}.
\end{quote}

\begin{quote}
\emph{anta maruttuvar ēḻaikaḷiṭam karuṇaiyuṭaṉ naṭantukoṇṭār.}
\end{quote}

\begin{quote}
That doctor treated the poor with compassion.
\end{quote}

-\ta{உடன்} \textquotedblleft{}with\textquotedblright{}: \ta{கருணையுடன்}, with compassion.

\begin{cousinweb}
\textbf{\ta{கருணை}} (\emph{karuṇai}) is borrowed from Sanskrit \emph{karuṇā}, \textquotedblleft{}pity\textquotedblright{} (CDIAL 2811); the final -ā becomes Tamil -\ta{ஐ}.
\end{cousinweb}

\begin{grammarlens}[title={What you've built}]
One more word for describing what you lived and felt.
\end{grammarlens}

\subsection*{Your turn}
\begin{itemize}
\raggedright
  \item \emph{Say it:} \emph{karuṇai}
  \item \emph{Say it:} \emph{karuṇai}, once more
  \item \emph{Say it:} \emph{anta maruttuvar ēḻaikaḷiṭam karuṇaiyuṭaṉ naṭantukoṇṭār.} — That doctor treated the poor with compassion.
  \item \emph{Recall:} say the Tamil for consolation, then the Tamil for to subside, then say \emph{karuṇai} again
\end{itemize}

\begin{tcolorbox}[breakable,colback=teal!4,colframe=teal!35!black,title={Before you move on}]
What does \ta{கருணை} mean? (\textquotedblleft{}Mercy; compassion\textquotedblright{}.) Say it once more.
\end{tcolorbox}
\section[\texorpdfstring{\ta{பரிவு} (parivu)}{பரிவு (parivu)}]{\ta{பரிவு} (parivu) — tenderness; affection}

\label{lesson:TA-C457-parivu}



\begin{quote}
Before the new one: say the Tamil for to subside, then the Tamil for mercy.
\end{quote}

\subsection*{You'll want to know: \ta{பரிவு}}
\textbf{\ta{பரிவு}} — \emph{parivu} — \textquotedblleft{}tenderness; affection\textquotedblright{}.

\begin{quote}
\ta{பாட்டி} \ta{பேரனிடம்} \ta{பரிவுடன்} \ta{பேசினாள்}.
\end{quote}

\begin{quote}
\emph{pāṭṭi pēraṉiṭam parivuṭaṉ pēciṉāḷ.}
\end{quote}

\begin{quote}
Grandmother spoke tenderly to her grandson.
\end{quote}

-\ta{இடம்} marks the one a feeling is shown to: \ta{பேரனிடம்}, to the grandson.

\begin{cousinweb}
\textbf{\ta{பரிவு}} (\emph{parivu}) is native Dravidian, a noun from \ta{பரி} \emph{pari}, \textquotedblleft{}to be affectionate\textquotedblright{} (DEDR 3964).
\end{cousinweb}

\begin{grammarlens}[title={What you've built}]
One more word for describing what you lived and felt.
\end{grammarlens}

\subsection*{Your turn}
\begin{itemize}
\raggedright
  \item \emph{Say it:} \emph{parivu}
  \item \emph{Say it:} \emph{parivu}, once more
  \item \emph{Say it:} \emph{pāṭṭi pēraṉiṭam parivuṭaṉ pēciṉāḷ.} — Grandmother spoke tenderly to her grandson.
  \item \emph{Recall:} say the Tamil for to subside, then the Tamil for mercy, then say \emph{parivu} again
\end{itemize}

\begin{tcolorbox}[breakable,colback=teal!4,colframe=teal!35!black,title={Before you move on}]
What does \ta{பரிவு} mean? (\textquotedblleft{}Tenderness; affection\textquotedblright{}.) Say it once more.
\end{tcolorbox}
\section[\texorpdfstring{\ta{அனுதாபம்} (aṉutāpam)}{அனுதாபம் (aṉutāpam)}]{\ta{அனுதாபம்} (aṉutāpam) — sympathy}

\label{lesson:TA-C457-anutapam}



\begin{quote}
Before the new one: say the Tamil for mercy, then the Tamil for tenderness.
\end{quote}

\subsection*{You'll want to know: \ta{அனுதாபம்}}
\textbf{\ta{அனுதாபம்}} — \emph{aṉutāpam} — \textquotedblleft{}sympathy\textquotedblright{}.

\begin{quote}
\ta{அந்தக்} \ta{குடும்பத்தின்} \ta{இழப்புக்கு} \ta{நாங்கள்} \ta{அனுதாபம்} \ta{தெரிவித்தோம்}.
\end{quote}

\begin{quote}
\emph{antak kuṭumpattiṉ iḻappukku nāṅkaḷ aṉutāpam terivittōm.}
\end{quote}

\begin{quote}
We expressed sympathy for that family’s loss.
\end{quote}

\ta{அனுதாபம்} \ta{தெரிவி}, to express sympathy.

\begin{cousinweb}
\textbf{\ta{அனுதாபம்}} (\emph{aṉutāpam}) is borrowed from Sanskrit \emph{anutāpa}, \textquotedblleft{}remorse; heat\textquotedblright{} (Monier-Williams); in Tamil it means feeling the pain of another.
\end{cousinweb}

\begin{grammarlens}[title={What you've built}]
One more word for describing what you lived and felt.
\end{grammarlens}

\subsection*{Your turn}
\begin{itemize}
\raggedright
  \item \emph{Say it:} \emph{aṉutāpam}
  \item \emph{Say it:} \emph{aṉutāpam}, once more
  \item \emph{Say it:} \emph{antak kuṭumpattiṉ iḻappukku nāṅkaḷ aṉutāpam terivittōm.} — We expressed sympathy for that family’s loss.
  \item \emph{Recall:} say the Tamil for mercy, then the Tamil for tenderness, then say \emph{aṉutāpam} again
\end{itemize}

\begin{tcolorbox}[breakable,colback=teal!4,colframe=teal!35!black,title={Before you move on}]
What does \ta{அனுதாபம்} mean? (\textquotedblleft{}Sympathy\textquotedblright{}.) Say it once more.
\end{tcolorbox}
\section[\texorpdfstring{\ta{பரிதாபம்} (paritāpam)}{பரிதாபம் (paritāpam)}]{\ta{பரிதாபம்} (paritāpam) — pitifulness; pity}

\label{lesson:TA-C457-paritapam}



\begin{quote}
Before the new one: say the Tamil for tenderness, then the Tamil for sympathy.
\end{quote}

\subsection*{You'll want to know: \ta{பரிதாபம்}}
\textbf{\ta{பரிதாபம்}} — \emph{paritāpam} — \textquotedblleft{}pitifulness; pity\textquotedblright{}.

\begin{quote}
\ta{மழையில்} \ta{நனைந்த} \ta{நாயைப்} \ta{பார்க்கப்} \ta{பரிதாபமாக} \ta{இருந்தது}.
\end{quote}

\begin{quote}
\emph{maḻaiyil naṉainta nāyaip pārkkap paritāpamāka iruntatu.}
\end{quote}

\begin{quote}
The dog soaked in the rain was a pitiful sight.
\end{quote}

\ta{பரிதாபமாக} \ta{இருக்கிறது}, it is a sorry sight.

\begin{cousinweb}
\textbf{\ta{பரிதாபம்}} (\emph{paritāpam}) is borrowed from Sanskrit \emph{paritāpa}, \textquotedblleft{}scorching heat; pain, grief\textquotedblright{} (Monier-Williams).
\end{cousinweb}

\begin{grammarlens}[title={What you've built}]
One more word for describing what you lived and felt.
\end{grammarlens}

\subsection*{Your turn}
\begin{itemize}
\raggedright
  \item \emph{Say it:} \emph{paritāpam}
  \item \emph{Say it:} \emph{paritāpam}, once more
  \item \emph{Say it:} \emph{maḻaiyil naṉainta nāyaip pārkkap paritāpamāka iruntatu.} — The dog soaked in the rain was a pitiful sight.
  \item \emph{Recall:} say the Tamil for tenderness, then the Tamil for sympathy, then say \emph{paritāpam} again
\end{itemize}

\begin{tcolorbox}[breakable,colback=teal!4,colframe=teal!35!black,title={Before you move on}]
What does \ta{பரிதாபம்} mean? (\textquotedblleft{}Pitifulness; pity\textquotedblright{}.) Say it once more.
\end{tcolorbox}
\section[\texorpdfstring{\ta{அக்கறை} (akkaṟai)}{அக்கறை (akkaṟai)}]{\ta{அக்கறை} (akkaṟai) — concern; care}

\label{lesson:TA-C457-akkarai}



\begin{quote}
Before the new one: say the Tamil for sympathy, then the Tamil for pitifulness.
\end{quote}

\subsection*{You'll want to know: \ta{அக்கறை}}
\textbf{\ta{அக்கறை}} — \emph{akkaṟai} — \textquotedblleft{}concern; care\textquotedblright{}.

\begin{quote}
\ta{அப்பா} \ta{என்} \ta{உடல்நலத்தில்} \ta{அக்கறை} \ta{காட்டினார்}.
\end{quote}

\begin{quote}
\emph{appā eṉ uṭalnalattil akkaṟai kāṭṭiṉār.}
\end{quote}

\begin{quote}
Father showed concern for my health.
\end{quote}

What one cares about takes -\ta{இல்}: \ta{படிப்பில்} \ta{அக்கறை}, concern for one’s studies.

\begin{cousinweb}
\textbf{\ta{அக்கறை}} (\emph{akkaṟai}) is native Dravidian (DEDR 21); in Kannada and Telugu the same word means need or necessity.
\end{cousinweb}

\begin{grammarlens}[title={What you've built}]
One more word for describing what you lived and felt.
\end{grammarlens}

\subsection*{Your turn}
\begin{itemize}
\raggedright
  \item \emph{Say it:} \emph{akkaṟai}
  \item \emph{Say it:} \emph{akkaṟai}, once more
  \item \emph{Say it:} \emph{appā eṉ uṭalnalattil akkaṟai kāṭṭiṉār.} — Father showed concern for my health.
  \item \emph{Recall:} say the Tamil for sympathy, then the Tamil for pitifulness, then say \emph{akkaṟai} again
\end{itemize}

\begin{tcolorbox}[breakable,colback=teal!4,colframe=teal!35!black,title={Before you move on}]
What does \ta{அக்கறை} mean? (\textquotedblleft{}Concern; care\textquotedblright{}.) Say it once more.
\end{tcolorbox}
